\section{Session: Memulai, Menyimpan, dan Menghapus Session}

\textbf{Session} adalah mekanisme server-side untuk menyimpan data pengguna selama interaksi dengan aplikasi web. Karena HTTP bersifat \textit{stateless} (setiap request independen), session menjembatani kebutuhan untuk ``mengingat'' pengguna antar halaman. Saat session dimulai, PHP menghasilkan \textbf{Session ID} unik yang dikirim ke browser sebagai cookie (biasanya bernama \code{PHPSESSID}). Pada request berikutnya, browser mengirimkan cookie ini kembali sehingga PHP dapat mengidentifikasi pengguna dan mengakses data session yang disimpan di server \cite{phpManual}.

Untuk menggunakan session, panggil \code{session\_start()} di awal skrip --- \textbf{sebelum} output apa pun (termasuk spasi atau baris kosong sebelum tag \code{<?php}). Data disimpan dalam superglobal \code{\$\_SESSION} sebagai array asosiatif. Menyimpan: \code{\$\_SESSION['user'] = 'cecep'}; membaca: \code{echo \$\_SESSION['user']}; menghapus satu key: \code{unset(\$\_SESSION['user'])}. Untuk mengakhiri session secara penuh (logout), gunakan kombinasi \code{session\_unset()} dan \code{session\_destroy()} \cite{phpRightWay}.

\begin{figure}[ht]
\centering
\begin{tikzpicture}[
  box/.style={draw, rounded corners, minimum width=3cm, minimum height=1cm, align=center, font=\small, fill=#1},
  arrow/.style={-Stealth, thick},
  node distance=1.2cm
]
  \node[box=blue!15] (start) {User mengakses\\halaman};
  \node[box=green!15, below=of start] (sess) {\code{session\_start()}};
  \node[box=orange!15, below=of sess] (check) {Session ID\\ada?};
  \node[box=yellow!15, below left=1.5cm and 2cm of check] (new) {Buat Session ID\\baru (server)};
  \node[box=yellow!15, below right=1.5cm and 2cm of check] (load) {Muat data\\session lama};
  \node[box=purple!15, below=2.5cm of check] (use) {Gunakan\\{\footnotesize\code{\$\_SESSION}}};
  \node[box=red!15, below=of use] (end) {Akhir skrip /\\\code{session\_destroy()}};

  \draw[arrow] (start) -- (sess);
  \draw[arrow] (sess) -- (check);
  \draw[arrow] (check) -- node[above left, font=\footnotesize]{Tidak} (new);
  \draw[arrow] (check) -- node[above right, font=\footnotesize]{Ya} (load);
  \draw[arrow] (new) |- (use);
  \draw[arrow] (load) |- (use);
  \draw[arrow] (use) -- (end);
\end{tikzpicture}
\caption{Flowchart siklus hidup session di PHP}
\end{figure}

Keamanan session sangat penting karena \textit{session hijacking} (pencurian Session ID) dapat memberikan akses tidak sah ke akun pengguna. Beberapa praktik keamanan yang direkomendasikan: (1) panggil \code{session\_regenerate\_id(true)} setelah login berhasil untuk mencegah \textit{session fixation}; (2) simpan Session ID di cookie dengan atribut \code{HttpOnly} dan \code{Secure}; (3) atur \code{session.cookie\_samesite = 'Strict'} untuk mencegah pengiriman cookie lintas situs; (4) batasi masa berlaku session; (5) periksa IP atau \textit{user-agent} untuk deteksi anomali. Konfigurasi ini dapat diatur melalui \code{php.ini} atau fungsi \code{session\_set\_cookie\_params()} \cite{owaspCsrf}.

\begin{webcode}[language=PHP, caption={Login sederhana dengan session dan keamanan}]
<?php
// login.php
session_start();

// Pengaturan keamanan cookie session
session_set_cookie_params([
  'lifetime' => 3600,        // 1 jam
  'path'     => '/',
  'secure'   => true,        // hanya HTTPS
  'httponly'  => true,        // tidak diakses JS
  'samesite' => 'Strict'     // cegah CSRF
]);

if ($_SERVER['REQUEST_METHOD'] === 'POST') {
  $user = htmlspecialchars($_POST['username'] ?? '');
  $pass = $_POST['password'] ?? '';

  // Simulasi: di produksi, cek dari database dengan password_verify()
  if ($user === 'admin' && $pass === 'rahasia123') {
    session_regenerate_id(true); // cegah session fixation
    $_SESSION['username']  = $user;
    $_SESSION['login_time'] = time();
    header("Location: dashboard.php");
    exit;
  } else {
    $error = "Username atau password salah.";
  }
}
?>
\end{webcode}

\begin{catatan}
Di lingkungan produksi, jangan pernah menyimpan password dalam teks biasa (\textit{plain text}). Gunakan \code{password\_hash(\$pass, PASSWORD\_DEFAULT)} untuk menyimpan dan \code{password\_verify(\$input, \$hash)} untuk memverifikasi. Session data disimpan di server (biasanya di folder \code{/tmp/}), sehingga lebih aman daripada cookie yang tersimpan di browser.
\end{catatan}

Untuk proses logout, tiga langkah diperlukan: (1) panggil \code{session\_start()} untuk mengakses session aktif; (2) hapus semua variabel session dengan \code{\$\_SESSION = []} atau \code{session\_unset()}; (3) hancurkan session di server dengan \code{session\_destroy()}. Opsional, hapus juga cookie session di browser dengan \code{setcookie(session\_name(), '', time()-3600)}. Langkah ini memastikan session benar-benar berakhir dan tidak dapat digunakan kembali \cite{phpManual}.

